\textbf{猕猴ECoG电极植入的长期稳定性评估} \hfill SCI期刊发表 \quad 2024.6 -- 2025.12 \titlebreak
项目简述：高密度µECoG能够以亚毫米分辨率探究神经功能，但受术中脑移位影响，电极定位存在数毫米偏差，阻碍电生理信号与功能图谱的精确配准。本研究通过透明颅窗实现在体视觉定位，并系统评估猕猴ECoG电极植入后的空间位置稳定性与电生理特性稳定性，为长期多模态神经接口研究提供技术基础。
\begin{itemize}[nosep]
    \item 针对脑移位导致的电极定位不确定性，应用基于血管交叉点基准的图像配准算法（透视变换+最小二乘拟合），量化术后14天内电极微位移量（中位0.20 mm）
    \item 建立ECoG信号预处理流程：数字滤波、降采样、共模重参考（CAR）及ICA伪迹剔除，确保信号质量
    \item 采用小波变换进行时频分解，提取gamma频段功率谱特征，频谱白化处理以补偿1/f衰减，增强可视化效果；构建MVPA解码框架：基于gamma功率时空特征（64通道×8时间点）构造高维特征向量，采用SVM结合10折交叉验证，实现红绿/黑白光栅刺激的高精度分类（>75\%准确率）并验证神经解码性能的稳定性
    \item 纵向评估5个月以上记录稳定性，验证组织覆盖状态下神经解码性能的持续有效性。
\end{itemize}
